% =============================================================================
% Project : THE_PHILOSOPHY_OF_VIABILITY
% Book    : Book I — Ontology of Viability
% Book ID : VS.B01
%
% File    : 4_ASSETS/diagrams/
%           VS.B01.P01.C03.S05.SS08.P032.DG01.components_composition_and_identity_relations.tex
%
% Asset   : DG01 — Components, Composition, and Identity Relations
%
% Purpose :
%   To represent the distinction between component plurality,
%   compositional organization, and identity without implying
%   that any one of these necessarily entails another.
%
% Requires (preamble):
%   \usepackage{tikz}
%   \usetikzlibrary{calc, positioning, fit}
% =============================================================================

\begin{figure}[htbp]
    \centering
    \hyphenpenalty=10000 \exhyphenpenalty=10000\relax  % no hyphenation inside diagram labels

    \begin{tikzpicture}[
        font=\small,
        node distance=0.4cm,
        % --- base box ---
        box/.style={
            draw, rounded corners=3pt, align=center,
            inner sep=6pt, line width=0.7pt
        },
        % --- node roles ---
        compbox/.style={box, text width=2.6cm, minimum height=1.1cm,
                        fill=blue!8, draw=blue!55!black},
        compositionbox/.style={box, text width=6.8cm, minimum height=1.5cm,
                               fill=blue!16, draw=blue!55!black},
        identitybox/.style={box, text width=6.8cm, minimum height=1.5cm,
                            dashed, fill=gray!10, draw=black!60},
        % --- frames and labels ---
        vsgroup/.style={draw=black!45, dashed, rounded corners=5pt,
                        inner sep=8pt, line width=0.5pt},
        grouplabel/.style={font=\footnotesize\itshape, text=black!75,
                           align=left},
        legendtext/.style={font=\scriptsize},
        % --- undirected lines (no arrowheads: no entailment, no priority) ---
        vsrel/.style={draw=black!70, line width=0.8pt},
        vslink/.style={draw=black!55, dashed, line width=0.7pt}
    ]

    % ---------------------------------------------------------------
    % Components (equal rank)
    % ---------------------------------------------------------------
    \node[compbox] (b) {Component B};
    \node[compbox, left=of b]  (a) {Component A};
    \node[compbox, right=of b] (c) {Component C};

    \node[vsgroup, fit=(a)(b)(c)] (comps) {};
    \node[grouplabel, anchor=south west] at (comps.north west) {
        Component plurality
    };

    % ---------------------------------------------------------------
    % Composition and identity
    % ---------------------------------------------------------------
    \node[compositionbox, below=1.1cm of comps] (whole) {
        \textbf{Composition}\\[3pt]
        Organized relations among components
    };
    \node[identitybox, below=1.1cm of whole] (id) {
        \textbf{Identity}\\[3pt]
        Unity or plurality requires identity-level grounds
    };

    % each component participates in the composition (straight, vertical)
    \draw[vsrel] (a.south) -- (a.south |- whole.north);
    \draw[vsrel] (b.south) -- (b.south |- whole.north);
    \draw[vsrel] (c.south) -- (c.south |- whole.north);

    % possible ontological relevance (not entailment)
    \draw[vslink] (whole) -- (id);

    % ---------------------------------------------------------------
    % Legend and interpretive constraint
    % ---------------------------------------------------------------
    \node[below=0.7cm of id, legendtext] (legend) {
        \tikz[baseline=-0.5ex]\draw[vsrel] (0,0) -- (0.6,0);~participation in an organized configuration
        \quad\quad
        \tikz[baseline=-0.5ex]\draw[vslink] (0,0) -- (0.6,0);~possible ontological relevance
    };

    \node[below=0.4cm of legend, text width=10.5cm, align=center,
          font=\footnotesize\itshape] {
        \textbf{Interpretive constraint:}
        The component-to-composition connections represent
        participation in an organized configuration. The
        composition-to-identity connection represents possible
        ontological relevance, not entailment. Component count
        and compositional unity do not determine identity number.
    };

    \end{tikzpicture}

    \caption{The distinct relations among components, composition, and identity}
    \label{fig:components-composition-and-identity-relations}
\end{figure}

% =============================================================================
% END OF ASSET DG01
% =============================================================================
